% Chapter 5 -- The Recall Bottleneck
% Standalone, compilable draft (FPT-template article class + booktabs), sibling to
% sec_4_results.tex. Root-causes the recall shortfall established in Chapter 4 and
% reports the failed repair strategies as a negative-results contribution.
%
% MERGE/NUMBERING NOTE: this file opens with \section{The Recall Bottleneck} for
% standalone compilation and numbering, but at final assembly the assembler
% promotes it to a chapter (Chapter 5). Forward/back references to sibling chapters
% and their tables use typed numerals ("Chapter~\ref{ch:detection-eval}", "Section~\ref{sec:results-ablations}", "Table~4.1")
% because the chapter files do not share a label namespace when built separately;
% Chapter-4 sub-tables are referenced by section + Finding number rather than by
% hard-coded table number, to stay robust to final renumbering.
%
% Honesty-phrasing rules applied (personal_agenda P1):
%   - recall denominator restated / cross-referenced (>=10 surviving points, Sec 4.1)
%   - "mean IoU" always qualified as "of matched detections"
%   - negative results (#21/#24, temporal aggregation) reported AS FINDINGS, not
%     apologised for
%   - the retracted "~0.74 hard limit" is corrected explicitly (#24 Correction / #34)
%   - two measurement eras kept separate: the pre-promotion seq-00 repair search
%     (baseline F1 0.844; #21/#24/temporal) vs the promoted-config seq-08 fix (#34)
%   - every number traces to a finding / session log artifact (cited Finding~\#NN /
%     session_history.md)
%
\chapter{The Recall Bottleneck}
\label{ch:recall-bottleneck}

Chapter~\ref{ch:detection-eval} established two facts about the headline operating point: the pipeline is
precision-saturated (false positives are low and flat, about \num{0.7} per frame),
and it is recall-limited, missing roughly one car in four
($R=\num{0.730}$ on sequence~08). This chapter explains \emph{why} the recall is
limited. We show that the shortfall is a clustering failure rather than a
classification failure (Section~\ref{sec:recall-diagnosis}), root-cause it to
HDBSCAN splitting large, close cars into fragments (Section~\ref{sec:split-analysis}),
report a search of repair strategies that all fail to generalise
(Section~\ref{sec:repairs}), and describe the one intervention that did help ---
coarsening the clustering resolution --- which forces a correction of an earlier,
overstated ``hard ceiling'' claim (Section~\ref{sec:resolution-fix}). The negative
results are reported as findings, not apologised for: mapping which repairs fail,
and why, is itself a contribution, and it is what justifies the residual limitation
carried to Chapter~\ref{ch:discussion}. Unless stated otherwise, the recall denominator throughout is
the one defined in Section~\ref{sec:eval-protocol} --- ground-truth car instances with at least
\num{10} points surviving preprocessing.

\section{The shortfall is a clustering failure, not a classifier failure}
\label{sec:recall-diagnosis}

The recall ceiling is upstream of the learned classifier. The stage decomposition
of Section~\ref{sec:results-ablations} (Findings~\#47/\#49) makes this direct: geometric clustering with the
size/shape filters but \emph{no} learned classifier and \emph{no} track filter
already reaches recall \num{0.775}, which is \emph{higher} than the full pipeline's
\num{0.730}. Every learned stage that follows can only hold or reduce that recall:
the classifier discards clusters, and the centroid tracker only re-labels and
filters clusters that geometric clustering has already produced --- it emits a
detection in a frame only where a cluster exists in that frame, and never
interpolates a detection across a gap (\texttt{src/tracker.py};
\texttt{src/evaluate.py}), so neither stage can invent a detection where clustering
found nothing. Across the full sequence the two stages together trade roughly
\num{1600} true positives (\num{27051} geometric-only versus \num{25478} full;
Findings~\#47/\#49) for the removal of $\sim$\num{152000} false positives; this is a
distinct trade from the $+\num{1655}$ true positives that the resolution change of
Section~\ref{sec:resolution-fix} later recovers. Whatever caps recall is therefore
already present in the clusters that geometric clustering produces, before the
classifier sees them.

The geometric filters are not the cap either. Table-level ablation of the
size/shape gates (Finding~\#22, sequence~00, \num{100} frames) shows that of the
\num{16454} HDBSCAN clusters in the sample, only \num{2651} (\SI{16.1}{\percent})
pass the filters, and \num{1827} of the rejected clusters do overlap a ground-truth
car --- dominated by \texttt{min\_volume} (\SI{67.9}{\percent} of the GT-matching
rejections) and \texttt{min\_points} (\SI{26.4}{\percent}); sequence~08 shows the
same pattern (\SI{66.5}{\percent} / \SI{28.0}{\percent}). The dominant gates are
\texttt{min\_volume} and \texttt{min\_points} --- the two that reject \emph{small}
clusters. This is consistent with the rejected clusters being the smaller
sub-fragments of cars that HDBSCAN has already split (the dominant fragment passing
while the splinters fail the size gates), which the split analysis of
Section~\ref{sec:split-analysis} then confirms is the prevailing failure mode. We
state this as an inference rather than a direct measurement: the ablation records
only that a rejected cluster overlaps \emph{some} ground-truth car
(\texttt{src/analyze\_clustering.py}), not whether that same car also retains a
surviving cluster, so it establishes that the rejections are dominated by small
fragments, not that every one belongs to an otherwise-detected car. Either way the
diagnosis points one stage further upstream, at the clustering itself.

\section{Root cause: HDBSCAN splits large, close cars}
\label{sec:split-analysis}

To quantify how clustering caps recall we measured, for every eligible ground-truth
car, how its points distribute across HDBSCAN clusters (Finding~\#23). Each car is
assigned to exactly one of four states that partition the population: \emph{ok} (one
dominant cluster, under \SI{30}{\percent} noise), \emph{split} (points spread over
two or more clusters), \emph{noisy} (one cluster but at or above \SI{30}{\percent}
noise), or \emph{all-noise}. A separate, cross-cutting flag marks \emph{merged} cars
(a car whose dominant cluster is shared with another car); a merged car still has a
dominant cluster, so it is \emph{not} a fifth partition class and is reported below
as a memo count rather than a table row that would be summed.
Table~\ref{tab:split} gives the result on both sequences.

\begin{table}[H]
  \centering
  \caption{Distribution of ground-truth cars across HDBSCAN clusters
    (Finding~\#23; pre-promotion configuration; sequence~00, \num{100} frames;
    sequence~08, held-out sample). The four states partition the population and sum
    to the denominator; \emph{split} is the dominant failure mode and \emph{merging}
    (the cross-cutting memo row, not summed) is negligible. Split cars are uneven ---
    the largest fragment holds a median \num{0.92} of the points on sequence~00 and
    \num{0.77} on sequence~08 --- but a large point-fraction does not guarantee a
    match (see text): recall shows most split cars are lost.}
  \label{tab:split}
  \begin{tabular}{lcccc}
    \toprule
    & \multicolumn{2}{c}{Seq 00 (\num{1631} cars)} & \multicolumn{2}{c}{Seq 08 (\num{811} cars)} \\
    \cmidrule(lr){2-3}\cmidrule(lr){4-5}
    Status & Count & \% & Count & \% \\
    \midrule
    ok (single dominant cluster) & 1117 & 68.5 & 509 & 62.8 \\
    split (2+ clusters)          & 514  & 31.5 & 301 & 37.1 \\
    noisy / all-noise            & 0    & 0.0  & 1   & 0.1  \\
    \midrule
    partition total              & 1631 & 100  & 811 & 100  \\
    \addlinespace
    \textit{merged} (cross-cutting) & 8 & 0.5 & 0 & 0.0 \\
    \bottomrule
  \end{tabular}
\end{table}

Three things follow. First, \textbf{splitting, not merging, is the failure mode}:
\SIrange{31}{37}{\percent} of cars split, while merging is at most
\SI{0.5}{\percent}. The pipeline almost never fuses two cars into one detection; it
repeatedly breaks one car into several. Second, \textbf{the split is worst where
cars are largest and nearest}: the split rate falls monotonically with range ---
\SI{66}{\percent} at \SIrange{0}{10}{\metre}, \SI{40}{\percent} at
\SIrange{10}{20}{\metre}, \SI{19}{\percent} at \SIrange{20}{30}{\metre}, and
\SI{4}{\percent} at \SIrange{30}{50}{\metre} (Finding~\#24). Large, close cars return many points, and
HDBSCAN finds internal density variations within that dense return --- the roof line
against the body, the near side against the far side --- and cuts the car along
them. This is the mechanism behind the counter-intuitive near-range recall dip
reported in Section~\ref{sec:results-distance} (Finding~\#50): the nearest bin has \emph{lower} recall than
the \SIrange{10}{20}{\metre} bin despite carrying far more points, because a close
car dissolves into fragments that no longer match the whole instance. Third,
\textbf{the cleanly-clustered fraction tracks recall, and splitting usually loses
the car}. The fraction of cars that form one dominant cluster
(\SI{68.5}{\percent} on sequence~00, \SI{62.8}{\percent} on sequence~08) sits just
\emph{below} the recall the pipeline achieves on the same configuration
(\num{0.739} and \num{0.699} respectively). Two things follow. The pipeline recovers
essentially every cleanly-clustered car, so that fraction is close to a floor on
recall --- not a ceiling; recall exceeds it only by the few split cars whose dominant
fragment happens to survive matching and filtering. And because recall
(\num{0.699}) barely exceeds the clean fraction (\SI{62.8}{\percent}) while
\SI{37.1}{\percent} of cars split, the great majority of split cars are \emph{not}
recovered: a large dominant point-fraction (median \num{0.77}--\num{0.92}) does not
translate into a match, because matching depends on spatial overlap and on the
fragment clearing the size filters, not on point count. Splitting is therefore the
mechanism by which cars are lost, which is exactly why it is the recall bottleneck.
This also predicts the resolution fix of Section~\ref{sec:resolution-fix}: coarsening
the voxel converts split cars into clean ones, which should lift recall by shrinking
the reliance on recovering splits --- and it does, narrowing the sequence-08 gap
between the clean fraction and recall from \num{7.1} points (\SI{62.8}{\percent}
versus \num{0.699}) to \num{3.1} points (\SI{69.9}{\percent} versus \num{0.730}).

Figure~\ref{fig:failure-zooms} shows the mechanism directly on three sequence~08
frames. Missed cars (orange) appear as broken rows of points close to the ego
vehicle, each row too small on its own to survive the size filters, while cleaner
mid- and far-range cars (green) are detected as single clusters.

\begin{figure}[htbp]
  \centering
  \includegraphics[width=\linewidth]{seq08_failure_zooms.png}
  \caption{Failure close-ups on sequence~08 (true positives green, false positives
    red, missed ground-truth cars orange; dashed outlines are ground-truth car
    footprints). Near-ego missed cars are fragmented into several small point groups
    rather than one cluster, matching the split mechanism of
    Table~\ref{tab:split}. Although the panel title mentions both merge and split,
    merging is measured at $\leq\!\SI{0.5}{\percent}$ (Table~\ref{tab:split}), so the
    misses shown are fragmentation, not fusion of neighbouring cars.}
  \label{fig:failure-zooms}
\end{figure}

\section{Repair strategies fail to generalise}
\label{sec:repairs}

Given that splitting caps recall, the obvious responses fall into two families:
change how clusters \emph{form} (a larger minimum cluster size, a different
clusterer, distance-adaptive parameters, or pre-clustering point accumulation) or
\emph{reassemble} fragments after HDBSCAN has split them. We tried six in total;
Table~\ref{tab:strategies} collects the four with a reproducible quantitative
outcome, and two further strategies (temporal accumulation and threshold lowering)
are discussed in prose below. All were evaluated against the same pre-promotion
sequence-00 baseline (full pipeline, classifier and track filter on;
$F_1=\num{0.844}$), and the load-bearing ones were re-checked on the held-out
sequence~08.

\begin{table}[H]
  \centering
  \caption{Recall-repair strategies against the pre-promotion baseline (full
    pipeline, sequence~00, \num{100} frames; held-out sequence-08 $F_1$ where
    measured). Every strategy either loses $F_1$ outright or fails to hold its
    sequence-00 gain on the held-out sequence. This is the \emph{pre-promotion}
    repair search --- it predates the coarse-voxel configuration of
    Section~\ref{sec:resolution-fix}, which is why the baseline here is
    $F_1=\num{0.844}$, not the promoted headline. Sources in the ``Src'' column are
    finding numbers. Two further strategies --- temporal point accumulation and
    lowering the size thresholds --- are discussed in the text but omitted here:
    temporal accumulation is single-sourced to a reverted feature in the session
    log (not reproducible under the research freeze), and threshold lowering
    produced no change in true positives; both belong in prose rather than the main
    results table.}
  \label{tab:strategies}
  \setlength{\tabcolsep}{4pt}
  \begin{tabular}{llcccc}
    \toprule
    Strategy & Src & Seq 00 $P$ & Seq 00 $R$ & Seq 00 $F_1$ & Seq 08 $F_1$ \\
    \midrule
    Baseline (HDBSCAN, MCS \num{10}) & \#24 & 0.984 & 0.739 & \textbf{0.844} & \textbf{0.829} \\
    BEV clustering (res \SI{0.30}{\metre})   & \#21 & 0.993 & 0.640 & 0.779 & --- \\
    Higher MCS \num{20}                       & \#24 & 0.976 & 0.755 & 0.852 & 0.801 \\
    Fragment merge (\SI{1.5}{\metre}, \num{30}\,pt) & \#24 & 0.931 & 0.755 & 0.834 & 0.815 \\
    Adaptive HDBSCAN (\num{30}/\num{15}/\num{10}) & \#24 & 0.935 & 0.748 & 0.831 & --- \\
    \bottomrule
  \end{tabular}
\end{table}

Each strategy fails for a structural reason, and the reasons are consistent with the
diagnosis of Section~\ref{sec:split-analysis}:

\begin{itemize}
  \item \textbf{Higher minimum cluster size (MCS \num{20})} is the most instructive.
    It is the only variant to \emph{improve} $F_1$ on sequence~00
    (\num{0.844}~$\rightarrow$~\num{0.852}), by discouraging the small splinter
    clusters. But it does not generalise: on the held-out sequence~08 it drops to
    $F_1$~\num{0.801}, \emph{below} the \num{0.829} baseline, because a larger
    minimum also erases genuinely sparse distant cars. The sequence-00 gain was
    over-fitting.
  \item \textbf{Post-clustering fragment merge} (join clusters within
    \SI{1.5}{\metre}) raises sequence-00 recall to \num{0.755} but loses precision
    (\num{0.984}~$\rightarrow$~\num{0.931}), because the same proximity rule also
    fuses nearby non-car structure --- walls, poles --- into oversized clusters. On
    held-out data the precision loss cancels the recall gain
    ($F_1$~\num{0.815}~$<$~\num{0.829}).
  \item \textbf{Distance-adaptive HDBSCAN} (a larger minimum cluster size near the
    ego, smaller far away) introduces boundary artefacts at the ring transitions and
    is worse than the global setting.
  \item \textbf{Bird's-eye-view clustering} (Finding~\#21), a cheaper
    connected-component alternative, collapses recall to \num{0.640}: projecting to
    two dimensions merges objects that overlap in the ground plane but are separated
    in height, and morphological cleanup destroys sparse distant clusters.
  \item \textbf{Temporal aggregation} --- accumulating object points over three
    consecutive frames before clustering, to densify sparse cars --- collapses the
    pipeline ($F_1$~\num{0.073}, recall \num{0.039}).\footnote{Source:
    \texttt{docs/session\_history.md} (2026-06 exploratory session). The
    \texttt{temporal\_window} feature was implemented and then fully reverted, so
    this result is on the pre-promotion baseline and was not re-run under the
    research freeze; it is reported here in prose, not in Table~\ref{tab:strategies},
    for that reason, and its value is the mechanism below rather than the exact
    figure.} Accumulating points does not densify
    the existing clusters; it reshapes the clustering landscape entirely
    (about \num{134000} accumulated points yield \num{645}+ clusters versus
    $\sim$\num{240} single-frame, of which fewer pass the geometric filters,
    \num{17} versus \num{34}).
  \item \textbf{Lowering the cluster-size thresholds} to admit sparser groups
    produced no change in true positives at all: the binding constraint is cluster
    quality, not the count threshold, so relaxing the threshold merely admits more
    noise (Finding~\#24; \texttt{docs/session\_history.md}).
\end{itemize}

The failure modes are strategy-specific --- over-fitting, precision loss, ring
artefacts, projection ambiguity, a reshaped clustering landscape --- but they share
a limit. The two strategies that reassemble after the split (fragment merge, and to
a lesser extent the proximity-based part of the adaptive scheme) run into the fact
that density-based clustering has no notion of objectness: once a dense car is
partitioned along an internal density gap, no proximity- or size-based rule can
reliably tell ``rejoin these two fragments of one car'' apart from ``do not join
these two adjacent cars.'' The strategies that instead change how clusters
\emph{form} --- a larger minimum size, a different clusterer, per-ring parameters,
pre-clustering accumulation --- avoid the reassembly trap but hit the same wall from
the other side: any single global density criterion that is loose enough to keep a
split car whole is also loose enough to fuse a close pair, so a sequence-00 gain
becomes a held-out loss. The one lever that helped, changing the clustering
\emph{resolution} (Section~\ref{sec:resolution-fix}), is also a formation change; it
works not by escaping this limit but by happening to sit at a resolution that closes
intra-car gaps before it starts closing inter-car ones. This is why the earlier
project conclusion --- that recall was fixed at a hard ceiling of about \num{0.74}
(Finding~\#24, original 2026-06-24 verdict) --- was reasonable given only the evidence
in Table~\ref{tab:strategies}, and also, as the next section shows, wrong: it read a
limit on \emph{these} repairs as a limit on the approach.

\section{What worked: clustering resolution, and the corrected ceiling}
\label{sec:resolution-fix}

The intervention that helped was not a repair at all but a change to the clustering
\emph{resolution}: re-voxelising the object cloud to a \SI{0.10}{\metre} grid
\emph{before} running HDBSCAN, rather than clustering the full-resolution
\SI{0.05}{\metre} cloud (Finding~\#34, adopted in the runtime optimisation of
Section~\ref{sec:results-runtime}). A coarser grid closes the small intra-car density gaps that HDBSCAN was
cutting along, so fewer splits form in the first place. It is a formation change like
the failed ones in Section~\ref{sec:repairs}, but it lands on the favourable side of
the trade described there --- coarse enough to close intra-car gaps, still fine
enough not to fuse neighbouring cars.

\begin{table}[H]
  \centering
  \caption{Effect of coarsening the clustering voxel from \SI{0.05}{\metre} to
    \SI{0.10}{\metre} on the split rate (Finding~\#34; sequence~08, \num{100}
    frames, \num{810} ground-truth cars). More cars cluster cleanly and none newly
    merge --- the coarsening removes splits without fusing neighbours.}
  \label{tab:splitrate}
  \begin{tabular}{lcc}
    \toprule
    & Full-res (\SI{0.05}{\metre}) & Coarse (\SI{0.10}{\metre}) \\
    \midrule
    clustered ``ok''            & 504 (62.2\%) & 566 (69.9\%) \\
    split                        & 305 (37.7\%) & 242 (29.9\%) \\
    merged                       & 0 (0.0\%)    & 0 (0.0\%) \\
    mean clusters per split car  & 3.3          & 2.7 \\
    \bottomrule
  \end{tabular}
\end{table}

Table~\ref{tab:splitrate} shows the split rate falling by \num{7.8} points
(\SI{37.7}{\percent}~$\rightarrow$~\SI{29.9}{\percent}) with \emph{zero} new
merges.\footnote{The full-resolution counts here (Finding~\#34: \num{810} cars, split
\SI{37.7}{\percent}) and those in Table~\ref{tab:split} (Finding~\#23: \num{811}
cars, split \SI{37.1}{\percent}) are separate runs of the same analysis on the same
pre-promotion configuration; they differ by at most one car and \num{0.6} percentage
points, and the percentages are unaffected.}
At the pipeline level, on the full sequence~08, this lifts recall from \num{0.699} to
\num{0.730} ($+\num{1655}$ true positives) with precision essentially unchanged
(\num{0.903}~$\rightarrow$~\num{0.905}) and $F_1$ from \num{0.788} to \num{0.808}
(Finding~\#34). The same coarsening also reduces clustering runtime, which is why it
was adopted; its detection benefit is the subject here and its runtime co-benefit is
reported in Section~\ref{sec:results-runtime}.

This result forces a correction to the earlier ceiling claim, which we state plainly
rather than bury (Finding~\#24, Correction of 2026-07-24). The ``$\sim$\num{0.74}
hard limit'' was \emph{partly a resolution artefact}, not a fundamental property of
the approach: the strategies of Section~\ref{sec:repairs} genuinely do fail, but none
of them tried this particular resolution setting, and changing it recovered part of
the loss. The standing recommendation is therefore
\texttt{cluster\_voxel\_size}~$=\SI{0.10}{\metre}$ in production, not ``accept
\num{0.74}.''

A smaller structural limit does remain, and it is architectural. Coarsening the voxel
further would eventually start merging adjacent cars, trading the split problem for a
merge problem; and the clustering-method benchmark of Section~\ref{sec:results-ablations} (Finding~\#43)
shows that no fixed-radius density clusterer, at any tested radius, reaches even
HDBSCAN's recall. The ceiling that remains is the one imposed by density-based
clustering having no model of what a car is: it groups points by proximity and
density alone, so it cannot be pushed to the recall of a learned object detector by
tuning. Reaching that level would require replacing the clustering stage with a
learned detector, which is outside the scope of this thesis (Chapter~\ref{ch:discussion}). The
contribution of this chapter is to have located the ceiling precisely --- it is in
the clustering stage, it is splitting rather than classification or merging, and it
is lower than first claimed --- and to have mapped which repairs cannot move it and
which one can.

\section{Summary}
\label{sec:recall-summary}

The recall shortfall is clustering-bound, not classifier-bound: the classifier and
track filter sit on top of a recall ceiling that geometric clustering already imposes
(Section~\ref{sec:recall-diagnosis}). The dominant mechanism is HDBSCAN splitting
large, close cars into fragments --- \SIrange{31}{37}{\percent} of cars, worst at
short range --- while merging is negligible, and the cleanly-clustered fraction sits
just below the achieved recall, so most split cars are lost
(Section~\ref{sec:split-analysis}). A search of six repair strategies --- higher
minimum cluster size, fragment merging, adaptive and BEV clustering, temporal
accumulation, and threshold lowering --- all fail to generalise, for a shared reason:
density-based clustering has no notion of objectness, so any criterion loose enough to
keep a split car whole also fuses adjacent cars, and reassembling fragments after the
split faces the same ambiguity (Section~\ref{sec:repairs}). The one effective lever was clustering
resolution: a coarser pre-clustering voxel prevents splits from forming, lifting
recall \num{0.699}~$\rightarrow$~\num{0.730} and correcting the earlier overstated
``hard ceiling'' (Section~\ref{sec:resolution-fix}). The residual limit is
architectural --- density clustering has no notion of objectness --- and is carried
to the limitations and future-work discussion of Chapter~\ref{ch:discussion}.
